% !TeX root = ../main.tex
% !TeX encoding = utf8

\chapter{Presentación de la organización}

\section{Campo de actividad}

CaixaBank es una entidad de banca universal con sede central en Valencia. Opera
principalmente en España y Portugal, ofreciendo una cartera amplia de productos y
servicios financieros dirigidos tanto a particulares como a empresas, instituciones y
grandes patrimonios. Sus líneas de negocio principales son las siguientes:

\begin{itemize}
  \item \textbf{Banca minorista:} cuentas corrientes y de ahorro, depósitos, créditos
    al consumo, hipotecas y medios de pago para particulares y autónomos.
  \item \textbf{Banca de empresas y corporativa:} financiación, tesorería, comercio
    exterior y asesoramiento financiero para pymes, grandes empresas e instituciones.
  \item \textbf{Banca privada y gestión de patrimonio:} asesoramiento personalizado,
    fondos de inversión y planes de pensiones para clientes de alto patrimonio.
  \item \textbf{Seguros:} a través de VidaCaixa (seguros de vida y pensiones) y
    SegurCaixa Adeslas (seguros de salud y generales), entre otras filiales aseguradoras.
  \item \textbf{Banca digital:} la aplicación CaixaBankNow, con más de 11 millones de
    usuarios activos digitales en España, constituye uno de los principales canales de
    relación con los clientes [\ref{ref:informe2024}].
\end{itemize}

En Portugal, CaixaBank opera a través de BPI (Banco BPI, S.A.), del que es accionista
mayoritario con un 100\% del capital desde 2019 [\ref{ref:informe2024}].

\section{Breve historia}

Los orígenes de CaixaBank se remontan a 1904, cuando se fundó en Barcelona la
\textit{Caixa de Pensions per a la Vellesa i d'Estalvis} (Caja de Pensiones para la Vejez
y el Ahorro), entidad orientada a la acción social y al ahorro popular. En 1990, esta
institución se fusionó con la \textit{Caja de Ahorros y Monte de Piedad de Barcelona}
para constituir \textit{la Caixa}, la mayor caja de ahorros de España durante décadas
[\ref{ref:fundacion}].

En el marco de la reestructuración del sistema financiero español iniciada tras la crisis
de 2008, la Caixa de Pensions procedió a la segregación de su actividad financiera.
En 2011, CaixaBank S.A. quedó constituida como sociedad anónima bancaria cotizada en la
Bolsa de Madrid, convirtiéndose en el vehículo a través del cual el grupo ejercería su
actividad bancaria [\ref{ref:cnmv}].

El hito más reciente y determinante de su historia fue la fusión con Bankia, completada en
marzo de 2021. Esta operación, la mayor integración bancaria de la historia reciente de
España, situó a CaixaBank como primer banco del país por cuota de mercado en banca
minorista, con una posición de liderazgo en depósitos, créditos, seguros de vida y fondos
de inversión [\ref{ref:informe2024}].

\section{Tamaño actual}

A cierre del ejercicio 2024, CaixaBank contaba con los siguientes datos de dimensión
[\ref{ref:informe2024}]:

\begin{itemize}
  \item \textbf{Empleados en España:} aproximadamente 40.000 personas.
  \item \textbf{Empleados en el grupo (incluyendo BPI):} más de 45.000 personas.
  \item \textbf{Red de sucursales en España:} más de 4.000 oficinas.
  \item \textbf{Clientes en España:} más de 20 millones.
  \item \textbf{Activos totales del grupo:} alrededor de 680.000 millones de euros.
\end{itemize}

Estas cifras consolidan a CaixaBank como la primera franquicia bancaria del mercado
español por número de clientes, red de distribución y cuota de mercado en los principales
productos de banca minorista.

\section{Forma jurídica y estructura de propiedad}

CaixaBank, S.A. es una \textbf{sociedad anónima} (S.A.) cotizada en el mercado continuo
de la Bolsa de Madrid e integrada en el índice selectivo IBEX-35. Como entidad de crédito,
está sujeta a la supervisión directa del Banco Central Europeo (BCE) en el marco del
Mecanismo Único de Supervisión (MUS), así como a la supervisión complementaria del Banco
de España [\ref{ref:cnmv}].

La estructura de propiedad al cierre de 2024 es la siguiente [\ref{ref:cnmv}]:

\begin{itemize}
  \item \textbf{CriteriaCaixa, S.A.U.:} accionista de control con aproximadamente el
    45\% del capital social. CriteriaCaixa es el vehículo inversor y gestor de activos de
    la \textbf{Fundación Bancaria ``la Caixa''}, institución de carácter social sin ánimo
    de lucro que constituye el sustrato fundacional del grupo.
  \item \textbf{BlackRock, Inc.:} inversor institucional con una participación
    significativa (en torno al 5\%).
  \item \textbf{Resto del capital:} distribuido entre inversores institucionales
    internacionales y accionistas minoritarios.
\end{itemize}

CaixaBank pertenece al \textbf{Grupo CaixaBank}, cuya cabecera es la propia
CaixaBank, S.A. El grupo integra filiales en banca, seguros, gestión de activos y
servicios auxiliares, entre las que destacan VidaCaixa, SegurCaixa Adeslas, CaixaBank
Asset Management, CaixaBank Payments \& Consumer y Banco BPI [\ref{ref:informe2024}].

\endinput
% -------------------------------------------------------------------
% FIN DEL CAPÍTULO 1
% -------------------------------------------------------------------
